\chapter{Change Request Workflow}
\label{ch:change-request-workflow}

This chapter walks through the complete lifecycle of an \texttt{ops.change-request} execution,
from the moment the runbook is started to change record closure. It includes a detailed example
of a database schema migration with a rollback path.

% ─────────────────────────────────────────────────────────────────────────────
\section{Lifecycle Overview}
\label{sec:cr-lifecycle}
% ─────────────────────────────────────────────────────────────────────────────

A change request execution has three phases:

\begin{description}
  \item[Pre-execution] The change-manager reviews and approves the request. Freeze window
    checks are performed. The change record is opened in the external system.
  \item[Execution] The DRI-guided steps run inside the \texttt{ops.change-request} wrapper.
    Evidence is captured at each step. Approval gates may pause execution.
  \item[Post-execution] The DRI signs off. The change record is closed. The evidence bundle
    is written and available for compliance review.
\end{description}

If any step in the execution phase fails, the rollback sequence is automatically initiated
before the run terminates.

% ─────────────────────────────────────────────────────────────────────────────
\section{Pre-Execution Phase}
\label{sec:cr-pre-execution}
% ─────────────────────────────────────────────────────────────────────────────

\subsection{Change-Manager Approval Gate}

When a runbook starts that contains an \texttt{ops.change-request} step, gert evaluates the
pre-execution requirements before running any steps. The change-manager from
\texttt{meta.roles.change-manager} must approve before execution begins.

The approval request is presented to the change-manager with:
\begin{itemize}
  \item The runbook name, kind, and description.
  \item All declared inputs (with sensitive values redacted per governance rules).
  \item The \texttt{change-id} and \texttt{risk} level.
  \item The names of steps within the wrapper and any rollback specification.
\end{itemize}

\subsection{Freeze Window Check}

If the gert server is configured with a change freeze calendar, gert checks whether the
current time falls within a freeze window before presenting the approval request. If a freeze
window is active:
\begin{itemize}
  \item Changes with \texttt{risk: low} or \texttt{risk: medium} are blocked.
  \item Changes with \texttt{risk: high} require an additional \texttt{ops.approval} gate
    explicitly authorizing emergency execution.
  \item Changes with \texttt{risk: emergency} bypass the freeze check.
\end{itemize}

% ─────────────────────────────────────────────────────────────────────────────
\section{Execution Phase}
\label{sec:cr-execution}
% ─────────────────────────────────────────────────────────────────────────────

After pre-execution approval is granted, the steps inside the \texttt{ops.change-request}
wrapper execute in declared order. The following rules apply during execution:

\begin{itemize}
  \item Steps execute sequentially. Parallel execution is not supported within a change
    request wrapper.
  \item The change request is considered \emph{open} from the moment the first inner step
    starts until either the wrapper completes successfully or the rollback completes.
  \item If a step fails (non-zero exit code, timeout, or explicit failure), gert pauses
    execution and presents the DRI with two options: \textbf{retry} (re-execute the failed
    step) or \textbf{rollback} (abandon the change and execute the rollback runbook).
  \item Evidence from steps within the wrapper is captured and tagged with the change request
    ID.
\end{itemize}

\subsection{Automatic Rollback Injection}

If the DRI chooses rollback (or if an unrecoverable error occurs), gert:
\begin{enumerate}
  \item Writes a \texttt{change/rollbackInitiated} trace event with the failing step ID and
    reason.
  \item Invokes the runbook specified in \texttt{rollback.runbook} with the inputs from
    \texttt{rollback.inputs}.
  \item Waits for the rollback runbook to complete.
  \item Writes a \texttt{change/rollbackCompleted} or \texttt{change/rollbackFailed} event.
  \item Closes the current run with outcome \texttt{rolled-back}.
\end{enumerate}

% ─────────────────────────────────────────────────────────────────────────────
\section{Post-Execution Phase}
\label{sec:cr-post-execution}
% ─────────────────────────────────────────────────────────────────────────────

After all inner steps complete successfully, the change request wrapper:
\begin{enumerate}
  \item Writes a \texttt{change/closed} trace event with the change ID, outcome, and DRI.
  \item If the gert server is configured with an external change management integration,
    it updates the ticket (e.g., ServiceNow, Jira) with the run ID, outcome, and evidence
    bundle URL.
  \item Returns control to the outer runbook. Steps after the \texttt{ops.change-request}
    wrapper execute normally (e.g., post-deploy health checks).
\end{enumerate}

% ─────────────────────────────────────────────────────────────────────────────
\section{Complete Example: Database Schema Migration}
\label{sec:cr-db-migration-example}
% ─────────────────────────────────────────────────────────────────────────────

The following is an annotated runbook for a database schema migration. It demonstrates the
full change request lifecycle, including pre-execution approval, DRI-guided execution,
evidence capture, and rollback specification.

\begin{minted}{yaml}
apiVersion: runbook/v2
kit: ops/v1

meta:
  name: migrate-user-schema
  kind: change
  description: >
    Apply Flyway migration V42__add_user_preferences.sql to the production
    PostgreSQL database. Requires change-board approval and DRI sign-off.
  inputs:
    change_ticket:
      from: prompt
      description: "Change ticket (e.g., CHG0042819)"
    db_host:
      from: prompt
      description: "Database host (e.g., postgres.prod.internal)"
    db_name:
      from: prompt
      description: "Database name (e.g., users_db)"
    migration_version:
      from: prompt
      description: "Flyway migration version (e.g., V42)"
  roles:
    dri: "@db-lead"
    approver: ["@change-board"]
    change-manager: "@change-board"
    observer: ["@audit-team"]
  sla:
    approval-timeout: 24h
    escalation-target: "@vp-engineering"
  governance:
    allowed_commands: [flyway, psql, pg_dump, kubectl]
    deny_env_vars: ["*_SECRET*", "*_PASSWORD", "PGPASSWORD"]
    redact:
      - pattern: "(password=)[^\\s]+"
        replace: "$1[REDACTED]"

flow:
  # ── Step 1: Back up the schema before migration ───────────────────────────
  - step:
      id: pre_migration_backup
      type: ops.cli
      title: "Take pre-migration schema backup"
      command: pg_dump
      args:
        - "--schema-only"
        - "-h"
        - "{{ .db_host }}"
        - "-d"
        - "{{ .db_name }}"
        - "-f"
        - "pre-migration-schema.sql"
      capture:
        dump_output: stderr
      evidence:
        auto: true
        label: "pre-migration schema dump"

  # ── Step 2: DRI reviews the migration script ──────────────────────────────
  - step:
      id: review_migration
      type: ops.manual
      title: "DRI reviews migration script"
      requires_role: dri
      instructions: |
        Review the migration script before proceeding:
          cat migrations/{{ .migration_version }}__add_user_preferences.sql

        Confirm:
        1. The migration is non-destructive (no DROP TABLE, no column removal).
        2. The migration is backward compatible with the current application.
        3. An undo migration script exists at:
           migrations/undo/{{ .migration_version }}__rollback_user_preferences.sql
      evidence:
        - type: ops.evidence.attestation
          label: "DRI migration review"
          required: true
          prompt: >
            I have reviewed migration {{ .migration_version }} and confirm it is
            non-destructive and backward compatible.

  # ── Step 3: Change board approval ────────────────────────────────────────
  - step:
      id: change_board_approval
      type: ops.approval
      title: "Change board approval for schema migration"
      description: |
        Database schema migration request:
          Database: {{ .db_name }} on {{ .db_host }}
          Migration: {{ .migration_version }}__add_user_preferences.sql
          Ticket: {{ .change_ticket }}
          Risk: medium (additive schema change, backward compatible)
          Rollback: Flyway undo migration available
      approvers: ["@change-board"]
      requires_any: true
      timeout: 24h
      on_timeout: escalate
      evidence:
        - type: ops.evidence.attestation
          label: "Change board migration approval"
          required: true
          prompt: >
            I approve migration {{ .migration_version }} for {{ .db_name }}
            under change ticket {{ .change_ticket }}.

  # ── Step 4: Execute migration with rollback protection ────────────────────
  - step:
      id: execute_migration
      type: ops.change-request
      title: "Apply Flyway migration {{ .migration_version }}"
      change-id: "{{ .change_ticket }}"
      risk: medium
      rollback:
        runbook: ./rollback-db-migration.runbook.yaml
        inputs:
          db_host: "{{ .db_host }}"
          db_name: "{{ .db_name }}"
          migration_version: "{{ .migration_version }}"
      steps:
        # 4a. Validate migration checksums
        - step:
            id: flyway_validate
            type: ops.cli
            title: "Validate Flyway migration checksums"
            command: flyway
            args:
              - validate
              - "-url=jdbc:postgresql://{{ .db_host }}/{{ .db_name }}"
              - "-locations=filesystem:migrations"
            capture:
              validate_output: stdout
            evidence:
              auto: true
              label: "Flyway validate output"

        # 4b. Run the migration
        - step:
            id: flyway_migrate
            type: ops.cli
            title: "Apply migration"
            command: flyway
            args:
              - migrate
              - "-url=jdbc:postgresql://{{ .db_host }}/{{ .db_name }}"
              - "-locations=filesystem:migrations"
              - "-target={{ .migration_version }}"
            capture:
              migrate_output: stdout
            evidence:
              auto: true
              label: "Flyway migrate output"

        # 4c. Verify migration status
        - step:
            id: flyway_info
            type: ops.cli
            title: "Verify applied migrations"
            command: flyway
            args:
              - info
              - "-url=jdbc:postgresql://{{ .db_host }}/{{ .db_name }}"
            capture:
              info_output: stdout
            evidence:
              auto: true
              label: "Flyway info post-migration"

  # ── Step 5: Post-migration verification ──────────────────────────────────
  - step:
      id: post_migration_check
      type: ops.manual
      title: "Verify application health after migration"
      requires_role: dri
      instructions: |
        1. Check application error logs: no DB-related errors in the last 5 minutes.
        2. Run the application's health endpoint:
           curl https://api.prod.internal/healthz/db
        3. Confirm the new columns exist:
           psql -h {{ .db_host }} -d {{ .db_name }} \
             -c "SELECT column_name FROM information_schema.columns \
                 WHERE table_name = 'users' AND column_name LIKE '%preferences%';"
        4. Take a screenshot of the health check output.
      evidence:
        - type: ops.evidence.screenshot
          label: "Application health post-migration"
          required: true
          description: "Screenshot of health endpoint and schema verification."
        - type: ops.evidence.attestation
          label: "DRI post-migration sign-off"
          required: true
          prompt: >
            I confirm migration {{ .migration_version }} was applied successfully.
            The application is healthy and the new columns are present.
\end{minted}
